
\subsubsection{Interpretation}

The results support a self-reading interpretation: when prompted to reason step-by-step, the model generates linguistic uncertainty markers that remain in-context when it subsequently generates a confidence estimate. The negative correlation (r = -0.315) indicates the model adjusts confidence downward when more hedging is present.

\subsubsection{Alternative Explanations}

\textbf{Difficulty confound.} Difficult questions may produce both more hedging and lower confidence without hedging causally affecting confidence. This interpretation still supports using hedging as a calibration signal but does not demonstrate the self-reading mechanism.

\textbf{Prompt artifact.} The prompt structure might induce the correlation through some artifact. This is considered less likely given the effect size and consistency across 817 items.

\textbf{Token-count confound.} Without the token-padding control, it cannot be empirically ruled out that longer outputs drive the correlation independently of uncertainty content.

\subsubsection{Limitations}

\textbf{Single model.} Only GPT-3.5-turbo was tested. Generalization to other architectures requires replication.

\textbf{Correlational evidence.} The findings are correlational. An intervention study manipulating hedging markers would provide stronger causal evidence.

\textbf{ECE comparison not executed.} The primary ECE comparison (CoT+confidence vs. single interventions) was not executed with real API calls. The sub-hypothesis H-E1 verifying ECE computation infrastructure ran in mock mode; actual calibration improvement was not measured.

\textbf{Token-padding control not executed.} The planned control condition was not run. Length artifacts cannot be empirically ruled out.

\textbf{Single dataset.} Only TruthfulQA was evaluated. Cross-domain transfer to MMLU or other datasets is untested.

\subsubsection{Practical Implications}

For practitioners, CoT+confidence prompting provides a zero-shot method for obtaining confidence estimates that correlate with uncertainty signals, without model fine-tuning or access to internal logits. Full ECE validation remains future work.
